
\slajdid{05-s021}
\begin{frame}[label=05-s021]{Papir i olovka}
\gusttekst\setlength{\parskip}{2pt}
$a-b={?}$\\
\textit{Ako je $a=b$ rezultat $=0$\\Ako je $a>b$ klasično oduzimanje\\Ako je $a<b$ rezultat $=-(b-a)$ radimo klasično oduzimanje $b-a$ i rezultatu menjamo znak}

Mogli bi ovu proceduru da sprovedemo i u realizaciji digitalnog sistema pogotovo na primer ako su nam negativni brojevi znak i apsoluta vrednost.\\
Treba nam komparator ali je njega lako napraviti u binarnom brojnom sistemu.\\Napravili smo ga.
\begin{columns}[t,onlytextwidth]
\column{.48\textwidth}Kada smo radili sabiranje radili\\smo samo sa pozitivnim vrednostima\\međutim ako imamo i negativne brojeve\\nama trebaju sabiranja
\column{.48\textwidth}Kada razmatramo oduzimanje radili\\smo samo sa pozitivnim vrednostima\\međutim ako imamo i negativne brojeve\\nama trebaju oduzimanja
\end{columns}
\begin{center}\begin{tikzpicture}[x=1cm,y=.43cm,font=\fontsize{9}{10}\selectfont]
\node at (0,1) {$a,b\geq0$};\node at (7,1) {$a,b\geq0$};
\foreach \id/\y/\f in {a/0/{a+b},b/-1/{(-a)+b},c/-2/{a+(-b)},d/-3/{(-a)+(-b)}}{\node (L\id) at (0,\y) {$\f$};}
\foreach \id/\y/\f in {a/0/{a-b},b/-1/{(-a)-b},c/-2/{a-(-b)},d/-3/{(-a)-(-b)}}{\node (R\id) at (7,\y) {$\f$};}
\foreach \l/\r in {a/c,b/d,c/a,d/b}{\draw[DEplava,<->] (L\l.east)--(R\r.west);}
\end{tikzpicture}
\end{center}
\end{frame}
